\documentclass[aspectratio=169,11pt]{beamer}

\usetheme{Madrid}
\usecolortheme{default}
\usefonttheme{professionalfonts}
\setbeamertemplate{navigation symbols}{}
\setbeamertemplate{footline}[frame number]

\usepackage{amsmath, amssymb, booktabs}

\title{Migration, Local Labor Demand}
\subtitle{And Urban Labor Reallocation}
\author{Special Topic 4: Urban and Labor -- Week 5}
\date{}

\begin{document}

\begin{frame}
  \titlepage
\end{frame}

\begin{frame}{Central question}
\begin{itemize}
    \item How do workers, firms, wages, rents, and commuting flows adjust when local labor demand changes?
    \item Migration is one margin in a spatial labor-adjustment system.
    \item Local shocks move people, jobs, firms, rents, commute flows, and neighborhoods.
    \item The key empirical task is incidence: who can adjust, who gains, and who remains exposed?
\end{itemize}
\end{frame}

\begin{frame}{Where Week 5 fits}
\begin{itemize}
    \item Week 1: job-residence bundles define local labor opportunity.
    \item Week 2: housing and moving costs limit access to productive places.
    \item Week 3: segregation and neighborhoods create unequal job access.
    \item Week 4: safety and environment change the feasible set for work.
    \item Week 5: local shocks trigger dynamic reallocation across the whole system.
\end{itemize}
\end{frame}

\begin{frame}{Spatial labor-adjustment value}
\[
V_{ijrt}
=
w_{jt}(D_{\ell t})
- R_{rt}(D_{\ell t})
- \tau_{jrt}
+ A_{rt}
- M_{irt}
+ \varepsilon_{ijrt}
\]
\small
\begin{tabular}{p{0.20\linewidth} p{0.64\linewidth}}
\toprule
Term & Interpretation \\
\midrule
$w(D)$ & workplace wage response to local demand \\
$R(D)$ & rent or housing-cost capitalization \\
$\tau$ & commute cost linking residence and workplace \\
$A$ & amenities and disamenities \\
$M$ & moving cost and attachment to place \\
\bottomrule
\end{tabular}
\end{frame}

\begin{frame}{Adjustment margins}
\small
\begin{tabular}{p{0.23\linewidth} p{0.35\linewidth} p{0.26\linewidth}}
\toprule
Margin & What moves & Incidence question \\
\midrule
Wages & pay, earnings, job quality & do workers share gains? \\
Employment & jobs, unemployment, LFPR & who remains attached? \\
Migration & inflows and outflows & who can move? \\
Commuting & residence-workplace flows & who accesses without moving? \\
Rents & housing costs, prices & who captures surplus? \\
Firms & entry, exit, relocation & are job ladders rebuilt? \\
Composition & sector, skill, neighborhood mix & who is being compared? \\
\bottomrule
\end{tabular}
\end{frame}

\begin{frame}{Migration and commuting}
\begin{itemize}
    \item Migration changes residence and the whole job-residence bundle.
    \item Commuting changes access to workplaces without changing residence.
    \item A job boom can raise workplace employment with little resident population growth.
    \item Commuting openness can lower wage pressure and spread gains across places.
    \item Housing constraints can convert labor demand into rent growth instead of worker inflow.
\end{itemize}
\end{frame}

\begin{frame}{Incidence after a local shock}
\small
\begin{tabular}{p{0.24\linewidth} p{0.36\linewidth} p{0.25\linewidth}}
\toprule
Group & Possible gain & Possible loss \\
\midrule
Incumbent workers & higher wages, job access & rent pressure, scarring \\
Migrants & entry to opportunity & moving costs, crowding \\
Commuters & access to wage gains & time and transport costs \\
Landlords/homeowners & capitalization & local risk if decline \\
Firms & thicker labor market & higher labor and land costs \\
Local government & tax base & service pressure, fiscal decline \\
\bottomrule
\end{tabular}
\end{frame}

\begin{frame}{Persistence and hysteresis}
\begin{itemize}
    \item Shocks can scar workers through job loss, lost tenure, sectoral mismatch, and weaker networks.
    \item Local employment can recover while incumbent workers remain on lower earnings paths.
    \item Workers can recover by leaving even when the origin place declines.
    \item Place recovery and person recovery are different estimands.
    \item Dynamic designs must define incumbents before the shock.
\end{itemize}
\end{frame}

\begin{frame}{Places versus people}
\small
\begin{tabular}{p{0.30\linewidth} p{0.48\linewidth}}
\toprule
Object & Outcomes \\
\midrule
Place & employment, population, establishments, rents, tax base, commuting flows \\
Incumbent worker & earnings, employment, occupation, commute, moving cost, welfare \\
Migrant & destination job access, housing cost, local networks, earnings path \\
Commuter & workplace wage access, travel burden, residence-based welfare \\
Firm & hiring costs, entry, exit, productivity, location choice \\
Owner or landlord & capitalization, rent income, wealth exposure \\
\bottomrule
\end{tabular}
\end{frame}

\begin{frame}{Gentrification as labor restructuring}
\begin{itemize}
    \item Gentrification changes more than rents and demographics.
    \item It can shift local service demand, job composition, business turnover, and commute patterns.
    \item Incumbent residents may gain from new jobs, safer amenities, or ownership wealth.
    \item Incumbent renters may lose through rent pressure, network disruption, and longer commutes.
    \item The labor question is which jobs are created, destroyed, upgraded, or made inaccessible.
\end{itemize}
\end{frame}

\begin{frame}{Diversification as worker insurance}
\begin{itemize}
    \item Specialized places offer upside but expose workers to sector-specific shocks.
    \item Diversified places may provide alternative local job ladders after downturns.
    \item The insurance value depends on skill transferability, housing access, and local hiring networks.
    \item Diversity can proxy for size, amenities, human capital, or institutions.
    \item Good designs distinguish resilience from selection into already-strong cities.
\end{itemize}
\end{frame}

\begin{frame}{Methods and data}
\small
\begin{tabular}{p{0.26\linewidth} p{0.32\linewidth} p{0.26\linewidth}}
\toprule
Setting & Design & Data \\
\midrule
Local demand shocks & shift-share, IV, event study & ACS, QWI, industry shares \\
Migration incidence & predicted-destination IV & IRS flows, settlement shares \\
Commuting access & transport or travel-time shocks & LEHD, LODES, GTFS \\
Persistence & dynamic local-shock panels & matched workers, CZ panels \\
Gentrification & job-location and tract panels & business data, rents, ACS \\
Diversification & shock by industry mix & sector shares, worker outcomes \\
Equilibrium incidence & spatial calibration & wages, rents, flows, housing \\
\bottomrule
\end{tabular}
\end{frame}

\begin{frame}{What a good design shows}
\begin{itemize}
    \item The shock is separated from endogenous local trends.
    \item The geography matches the mechanism: residence, workplace, tract, county, CZ, or flow market.
    \item Wages are interpreted with rents, commuting, migration, and composition.
    \item Incumbents, migrants, commuters, firms, and landlords are not collapsed into one effect.
    \item The design states whether the target is nominal earnings, real access, resilience, or welfare.
\end{itemize}
\end{frame}

\begin{frame}{Week 5 lab}
\begin{itemize}
    \item Reproduce: synthetic local demand shock with wages, rents, population, commuting, and firms.
    \item Diagnose: classify outputs as place objects or person objects and assign incidence.
    \item Transfer: gentrification scenario with local service demand and rent pressure.
    \item Transfer: diversification scenario with specialized and diversified places hit by a sector shock.
    \item Goal: interpret an adjustment vector before interpreting one coefficient.
\end{itemize}
\end{frame}

\begin{frame}{Takeaways}
\begin{itemize}
    \item Migration and commuting are linked margins in local labor adjustment.
    \item Local shocks have incidence across workers, migrants, commuters, firms, and property owners.
    \item Place recovery is not the same thing as worker recovery.
    \item Gentrification and diversification are labor topics when they restructure jobs and risk.
    \item Week 6 turns these objects into research designs.
\end{itemize}
\end{frame}

\end{document}
